\clearpage
\section{Rechnungen zu den neuen Abschnitten}
\label{app:new-calcs}

Anhang~\ref{app:wolfram-output} belegt die Rechnungen der Fassung 0.1.0, Anhang~\ref{sec:anhang-b} die Herleitungen der
OP5-Fallstudie. Dieser Anhang belegt die Zahlen, die danach in Abschnitt~\ref{sec:one-number} und in Abschnitt~\ref{sec:open-points}
hinzugekommen sind. Die Rechnungen stammen aus den Companion-Runden v4.10.0-r52 bis v4.10.0-r54 und v4.10.0-r58 \cite{layers}. Skripte und vollständige
Originalausgaben liegen unter \texttt{paper/\allowbreak preprint-de/\allowbreak audit/\allowbreak anhang-c/}; die Fenster unten zeigen
die maßgeblichen Zeilen der Ausgaben unverändert. Drei Rechnungen liefen in Wolfram Language über den Wolfram-Kern der jeweiligen
Sitzung, eine in Python mit sympy und numpy, eine Kontrolle als Quantenschaltung mit Qiskit.

\newcommand{\neuausgabe}[3]{%
  \VerbatimInput[firstline=#2,lastline=#3,fontsize=\footnotesize,baselinestretch=.88,%
    breaklines,breakanywhere,frame=single]{audit/anhang-c/#1}}

\subsection{Zählregeln als Ähnlichkeitsgleichungen (Runde v4.10.0-r53)}
\label{app:new-counting}

Skript \texttt{op1\_\allowbreak 148\_\allowbreak moran\_\allowbreak overlap.wl}, Wolfram Language, 23.~November 2026. Die
Zählgleichung
\[
 \sum_j g(j)\,e^{-2\pi\gamma\sqrt{j(j+1)}}=1
\]
wird mit 30 Stellen gelöst, die Spinsumme bei \(j\le200\) abgeschnitten (Rest etwa \(10^{-130}\)). Die erste Gruppe gibt die Werte der Zählregeln und den Wert \(\gamma_1\) bei zusammenfallenden Extremzuständen,
die zweite Zeile prüft, dass mit den Gibbs-Gewichten Entropie durch Lyapunov-Exponent gleich \(\gamma\) ist.
\neuausgabe{op1_148_moran_overlap_stdout.txt}{5}{8}
\neuausgabe{op1_148_moran_overlap_stdout.txt}{12}{12}
Die zweite Gruppe zählt in einem abgeschnittenen System (Spins \(j\le1\)) die verschiedenen verketteten Abbildungen und gibt die
Entropierate durch den Lyapunov-Exponenten für Verkettungen der Länge \(n=1,\dots,4\): für verschiedene Kopien, für gespiegelte mit
eigenem Zentrum, für identische und für gespiegelte mit gemeinsamem Zentrum.
\neuausgabe{op1_148_moran_overlap_stdout.txt}{23}{26}
Ein früher Lauf dieser Gruppe hatte einen Konstruktionsfehler (gemeinsame Fixpunkte) und ist in der Ausgabedatei als verworfen vermerkt.

\subsection{Die Frame-Abbildung \texorpdfstring{\(J_{\mathrm B}\)}{JB} am Flächenoperator (Runde v4.10.0-r52)}
\label{app:new-jb}

Skript \texttt{k241\_\allowbreak h121\_\allowbreak jb\_\allowbreak area.py}, Python 3.12 mit sympy und numpy, 23.~November 2026;
Vorab-Notiz zu H121 vor der Rechnung festgeschrieben. Die Zeilen zeigen die Konjugation der Lorentz-Generatoren, die induzierte
Matrix \(\Gamma'\) bei \(\gamma_M\) und die größte Abweichung der Spektren unter Konjugation.
\neuausgabe{k241_out.txt}{5}{14}
\neuausgabe{k241_out.txt}{29}{29}

\subsection{Die Auslese als Extremum (Runde v4.10.0-r54, H122)}
\label{app:new-readout}

Skript \texttt{h122\_\allowbreak readout\_\allowbreak extremum.wl}, Wolfram Language, 23.~November 2026; Vorab-Notiz zu H122 vor der
Rechnung festgeschrieben. Die erste Gruppe gilt für das Feld mit parallelem elektrischem und magnetischem Teil: die Auslese eines
geboosteten Beobachters relativ zum Ruhesystem. Die zweite gilt für die Gezeitenauslese des Vergleichsmodells: eine einzelne Mode, das
Azimutmittel und das isotrope Mittel unter Boosts des Detektors entlang und quer zur Achse.
\neuausgabe{h122_readout_extremum.txt}{6}{8}
\neuausgabe{h122_readout_extremum.txt}{16}{16}
\neuausgabe{h122_readout_extremum.txt}{27}{27}
\neuausgabe{h122_readout_extremum.txt}{20}{22}
Zwei Läufe dieser Rechnung brachen mit einem Serverfehler des Wolfram-Konnektors ab und wurden ohne die Minimierungsteile wiederholt; die
Ausgabedatei vermerkt das. Die Zeichen in den Fenstern sind die der Skripte: \texttt{et} ist die Rapidität, \texttt{th} und \texttt{t}
sind Winkel zur Achse, \texttt{b} ist die Geschwindigkeit des Detektors.

\subsection{Die zeitlich geordnete Oktant-Auslese (Runde v4.10.0-r58, H123)}
\label{app:new-temporal}

Skript \texttt{h123\_\allowbreak temporal\_\allowbreak octant.wl}, Wolfram Language, 23.~November 2026; Vorab-Notiz zu H123 vor der Rechnung
festgeschrieben. Die erste Gruppe prüft exakt in 640 Fällen die Formel \(B_\tau=2^{-3j}e^{i\tau\varepsilon j\pi/2}\) für Spin \(j\le20\) und die
Wirkung der Zeitumkehr des Protokolls. Die zweite nennt die Spins, deren Orientierungen die Auslese trennt; die dritte gibt die Zählregel mit
Gewicht 1 auf den nicht getrennten Spins und ihren Abgleich mit den Datenfenstern.
\neuausgabe{h123_temporal_octant.txt}{13}{16}
\neuausgabe{h123_temporal_octant.txt}{59}{60}
\neuausgabe{h123_temporal_octant.txt}{67}{71}
Als unabhängige Kontrolle berechnet eine Quantenschaltung (Zyklustest mit Hilfsqubit, Qiskit, Skript \texttt{k244\_\allowbreak h123\_\allowbreak qiskit\_\allowbreak cycle\_\allowbreak test.py})
die Invariante für \(j=\tfrac12,1,2\); die größte Abweichung von der Formel ist \(6{,}6\cdot10^{-16}\).
\neuausgabe{k244_out.txt}{19}{19}
\neuausgabe{k244_out.txt}{36}{36}
\neuausgabe{k244_out.txt}{53}{55}
